\subsection{Cấu trúc xử lý trong một lượt hội thoại}

Trong một lượt hội thoại, message của người dùng lần lượt đi qua các nhóm NLU, tracker, policy, NLG và custom action. Các khối này liên kết với nhau theo đúng thứ tự thực thi trong Rasa thay vì hoạt động độc lập.

\begin{figure}[H]
    \centering
    \includegraphics[width=\textwidth,height=0.82\textheight,keepaspectratio]{\detokenize{content/source/rasa/sequence/seq_2_2_4_dialogue_turn.png}}
    \caption{Biểu đồ tuần tự cấu trúc xử lý trong một lượt hội thoại}
    \label{fig:rasa-dialogue-processing-groups}
\end{figure}

Hình~\ref{fig:rasa-dialogue-processing-groups} cho thấy các điểm quyết định quan trọng trong một lượt hội thoại. Sau khi \texttt{PhoBERTIntentClassifier} và các thành phần entity hoàn tất xử lý, \texttt{FallbackClassifier} kiểm tra confidence và độ nhập nhằng của intent. Nếu đạt ngưỡng, Rasa Core tiếp tục chọn action bằng \texttt{RulePolicy} và \texttt{MemoizationPolicy}; nếu không đạt, luồng sẽ chuyển sang nhánh fallback. Với các action dạng \texttt{utter\_*}, hệ thống lấy phản hồi từ \texttt{domain.yml}; còn với các action nghiệp vụ hoặc truy xuất tri thức, Action Server và LightRAG mới được gọi.

